% This is a report prepared with the LLNCS macro package for
% Springer Computer Science proceedings.
%
\documentclass[runningheads]{llncs}
%
\usepackage[T1]{fontenc}
% T1 fonts will be used to generate the final print and online PDFs,
% so please use T1 fonts in your manuscript whenever possible.
%
\usepackage[utf8]{inputenc}
\usepackage[spanish]{babel}
\usepackage{amsmath}
\usepackage{booktabs}
\usepackage{graphicx}
\usepackage{url}
% If you use the hyperref package, please uncomment the following two lines
% to display URLs in blue roman font according to Springer's eBook style:
%\usepackage{color}
%\renewcommand\UrlFont{\color{blue}\rmfamily}
%\urlstyle{rm}
%

\title{Desarrollo de una Aplicaci\'on para Optimizar Rutas de Reparto mediante un Algoritmo Gen\'etico}

\titlerunning{Aplicaci\'on VRP con Algoritmo Gen\'etico}

\author{Nombre del estudiante\inst{1}}
\authorrunning{Nombre del estudiante}
\institute{Instituci\'on educativa\\
\email{correo@ejemplo.com}}

\begin{document}

\maketitle              % typeset the header of the contribution

\begin{abstract}
Este reporte describe el proceso completo de creaci\'on de una aplicaci\'on para
resolver una instancia del Problema de Ruteo de Veh\'iculos (VRP, por sus siglas
en ingl\'es), aplicado a la distribuci\'on de paquetes desde un dep\'osito hacia un
conjunto de clientes. La soluci\'on desarrollada integra una interfaz gr\'afica de
escritorio, administraci\'on de datos mediante archivos CSV, selecci\'on de
destinos desde un mapa, visualizaci\'on geogr\'afica de rutas y un algoritmo
gen\'etico para buscar soluciones de baja distancia bajo restricciones de
capacidad y cantidad de veh\'iculos. El documento sigue una metodolog\'ia
cient\'ifica: primero contextualiza el problema log\'istico, despu\'es define los
objetivos del proyecto, presenta el dise\~no matem\'atico y computacional,
describe la experimentaci\'on realizada y finalmente eval\'ua los resultados,
limitaciones y oportunidades de mejora.

\keywords{VRP \and Algoritmo gen\'etico \and Optimizaci\'on \and Log\'istica \and Mapas interactivos}
\end{abstract}

\section*{Repositorio del proyecto}

El c\'odigo fuente, los archivos de datos y los recursos necesarios para ejecutar
la aplicaci\'on se encuentran disponibles en el siguiente repositorio de GitHub:
\url{https://github.com/usuario/nombre-del-repositorio}.

\section{Introducci\'on}

La log\'istica es una actividad fundamental para el funcionamiento de empresas,
instituciones y servicios que dependen del movimiento de mercanc\'ias. En un
sistema de distribuci\'on, no basta con contar con veh\'iculos y productos; tambi\'en
es necesario decidir qu\'e veh\'iculo atiende a cada cliente, en qu\'e orden se
realizan las entregas, cu\'anta carga transporta cada unidad y cu\'al es la
distancia total que se recorrer\'a. Estas decisiones influyen directamente en el
costo operativo, el tiempo de entrega, el consumo de combustible, el desgaste de
los veh\'iculos y la calidad del servicio ofrecido al cliente.

En contextos reales, la distribuci\'on de mercanc\'ias se vuelve m\'as compleja por
la presencia de restricciones. Un cami\'on no puede transportar carga ilimitada,
las entregas deben partir de un punto de origen, los clientes est\'an ubicados en
diferentes coordenadas geogr\'aficas y la distancia entre dos puntos no siempre
coincide con una l\'inea recta. Adem\'as, cuando aumenta el n\'umero de destinos,
crece r\'apidamente la cantidad de combinaciones posibles para visitar a todos los
clientes. Por esta raz\'on, la planeaci\'on manual de rutas puede producir
soluciones poco eficientes o requerir demasiado tiempo.

El problema estudiado en este proyecto se conoce como Problema de Ruteo de
Veh\'iculos o Vehicle Routing Problem (VRP). De manera general, el VRP consiste
en encontrar un conjunto de rutas para una flota de veh\'iculos que salen de un
dep\'osito, atienden a un conjunto de clientes y regresan al punto de origen. El
objetivo com\'un es minimizar la distancia total recorrida o el costo total de
transporte, respetando restricciones como capacidad de carga, cantidad de
veh\'iculos disponibles y demanda de cada cliente~\cite{toth2014}.

El VRP es un problema complejo porque pertenece a la familia de problemas de
optimizaci\'on combinatoria. Esto significa que la soluci\'on se busca entre una
gran cantidad de combinaciones discretas. Incluso una instancia peque\~na puede
generar muchas alternativas: primero debe decidirse el orden de visita de los
clientes y, adem\'as, c\'omo se dividen esos clientes entre los veh\'iculos. Si se
pretendiera revisar exhaustivamente todas las combinaciones, el tiempo de
c\'alculo crecer\'ia de forma acelerada conforme aumenta el n\'umero de clientes.

Debido a esta dificultad, en muchos escenarios se utilizan m\'etodos
aproximados. Estos m\'etodos no siempre garantizan encontrar el \'optimo global,
pero pueden obtener soluciones de buena calidad en tiempos razonables. Entre
ellos se encuentran las metaheur\'isticas, como b\'usqueda tab\'u, recocido
simulado, colonias de hormigas y algoritmos gen\'eticos. En este proyecto se
elige un algoritmo gen\'etico porque su funcionamiento se adapta bien a problemas
basados en permutaciones, como el orden de visita de clientes.

Los algoritmos gen\'eticos se inspiran en procesos evolutivos. Trabajan con una
poblaci\'on de soluciones candidatas, eval\'uan su calidad mediante una funci\'on
de aptitud, seleccionan soluciones prometedoras y generan nuevas alternativas
mediante operadores de cruza y mutaci\'on~\cite{goldberg1989}. Para el VRP, cada
soluci\'on puede representarse como una permutaci\'on de clientes. Despu\'es, esa
permutaci\'on se interpreta como una secuencia que se divide en rutas de acuerdo
con la capacidad de los veh\'iculos. Esta representaci\'on permite explorar muchos
\'ordenes de entrega sin construir manualmente cada ruta desde cero.

La aplicaci\'on desarrollada busca demostrar este enfoque de manera pr\'actica.
El sistema fue construido en Python y cuenta con una interfaz de escritorio en
Tkinter. Desde la aplicaci\'on, el usuario puede definir la cantidad de veh\'iculos,
la capacidad de carga, el dep\'osito, los destinos de entrega y los par\'ametros
del algoritmo gen\'etico. Los clientes pueden cargarse desde archivos CSV,
editarse en una tabla o agregarse mediante clics sobre un mapa. Una vez
ejecutada la optimizaci\'on, el sistema muestra el resumen por veh\'iculo, la ruta
asignada, la carga total, la distancia aproximada y un mapa interactivo con los
recorridos.

Una parte importante del proyecto es el c\'alculo de distancias. Para escenarios
locales, la aplicaci\'on usa OSMnx y NetworkX para obtener una red vial basada en
OpenStreetMap y calcular rutas sobre calles reales~\cite{osmnx,networkx}. Para
escenarios nacionales o con puntos demasiado separados, se implementa un modo
aproximado basado en la distancia Haversine multiplicada por un factor de
carretera. De este modo, la aplicaci\'on permite experimentar tanto con rutas
urbanas o regionales como con entregas a estados lejanos, aunque con diferente
nivel de precisi\'on.

El prop\'osito del proyecto es educativo y aplicado. Por un lado, permite
entender c\'omo se modela un problema log\'istico real como un problema de
optimizaci\'on. Por otro lado, muestra c\'omo integrar un algoritmo gen\'etico con
una interfaz gr\'afica, archivos de datos y mapas interactivos. La aplicaci\'on no
pretende sustituir a un sistema profesional de planeaci\'on log\'istica, pero s\'i
ofrece una base funcional para estudiar el VRP y visualizar sus resultados.

El resto del reporte se organiza de la siguiente manera. La Secci\'on 2 presenta
el objetivo general y los objetivos espec\'ificos del proyecto. La Secci\'on 3
describe el dise\~no del sistema, incluyendo el modelo matem\'atico, las
restricciones, la arquitectura, la interfaz, el algoritmo gen\'etico, el flujo de
ejecuci\'on y los diagramas de apoyo. La Secci\'on 4 explica la experimentaci\'on,
el ambiente de desarrollo, los datos utilizados, los par\'ametros y los resultados
observados. Finalmente, la Secci\'on 5 presenta las conclusiones, eval\'ua el
cumplimiento de los objetivos y propone mejoras futuras.

\section{Objetivos}

\subsection{Objetivo general}

Desarrollar una aplicaci\'on de escritorio capaz de optimizar rutas de reparto
para varios veh\'iculos mediante un algoritmo gen\'etico, considerando la distancia
recorrida, la capacidad m\'axima de los camiones y el n\'umero de veh\'iculos
disponibles.

\subsection{Objetivos espec\'ificos}

\begin{itemize}
    \item Modelar el problema de reparto como un VRP con dep\'osito, clientes,
    demandas, veh\'iculos y capacidad m\'axima. Este objetivo permite transformar
    una situaci\'on log\'istica real en un problema computacional medible, donde
    cada destino se representa como un nodo y cada carga como una demanda.
    \item Implementar un algoritmo gen\'etico con representaci\'on por
    permutaciones, cruza OX y mutaci\'on por intercambio. Este objetivo aporta el
    mecanismo de b\'usqueda de soluciones, evitando depender de una revisi\'on
    manual o exhaustiva de todas las combinaciones posibles.
    \item Incorporar penalizaciones cuando una soluci\'on exceda la capacidad de
    los veh\'iculos o requiera m\'as camiones que los disponibles. Con ello, el
    sistema puede evaluar soluciones no factibles y favorecer aquellas que
    cumplen las restricciones operativas.
    \item Crear una interfaz gr\'afica que permita cargar, agregar, modificar y
    eliminar destinos de entrega. Este objetivo facilita el uso del sistema por
    parte de usuarios que no necesitan interactuar directamente con el c\'odigo.
    \item Generar mapas interactivos para visualizar el dep\'osito, los clientes
    y las rutas resultantes. La visualizaci\'on permite interpretar r\'apidamente
    la soluci\'on y comprobar si las rutas son coherentes con la ubicaci\'on de los
    destinos.
    \item Evaluar el comportamiento de la aplicaci\'on con datos locales y
    nacionales aproximados. Esta evaluaci\'on permite observar el funcionamiento
    del sistema en escenarios de diferente escala y reconocer las limitaciones
    de cada modo de distancia.
\end{itemize}

\section{Dise\~no}

\subsection{Descripci\'on detallada del problema VRP}

El sistema resuelve una variante capacitada del VRP. Se cuenta con un dep\'osito
central y un conjunto de clientes que deben ser atendidos por una flota limitada
de veh\'iculos. Cada cliente tiene una demanda, entendida como peso, volumen o
unidad de carga. Cada veh\'iculo tiene una capacidad m\'axima, por lo que no puede
transportar una suma de demandas mayor que dicho l\'imite. Adem\'as, todos los
veh\'iculos disponibles se consideran homog\'eneos, es decir, tienen la misma
capacidad y el mismo criterio de costo asociado a la distancia.

La dificultad del problema surge de dos decisiones simult\'aneas. La primera es
la asignaci\'on: determinar qu\'e clientes ser\'an atendidos por cada veh\'iculo. La
segunda es el secuenciamiento: decidir el orden en que cada veh\'iculo visitar\'a
sus clientes. Una mala asignaci\'on puede producir rutas sobrecargadas o exigir
m\'as camiones de los disponibles. Un mal secuenciamiento puede aumentar
innecesariamente la distancia recorrida aunque la asignaci\'on sea factible.

En la aplicaci\'on, la soluci\'on se expresa como rutas cerradas. Cada ruta inicia
en el dep\'osito, visita uno o m\'as clientes y regresa al dep\'osito. Este
comportamiento se refleja tanto en la tabla de resultados como en el mapa
generado. El dep\'osito se identifica visualmente como origen y destino final de
todos los veh\'iculos.

\subsection{Modelo matem\'atico utilizado}

Sea $V = \{0,1,2,\ldots,n\}$ el conjunto de nodos, donde el nodo $0$ representa
el dep\'osito y los nodos $1$ a $n$ representan clientes. Cada cliente $i$ tiene
una demanda $d_i$, y el dep\'osito tiene demanda $d_0=0$. Sea $M$ el n\'umero de
veh\'iculos disponibles y $Q$ la capacidad m\'axima de cada veh\'iculo. La distancia
entre dos nodos $i$ y $j$ se denota como $c_{ij}$.

La versi\'on cl\'asica del VRP puede expresarse con variables binarias $x_{ijk}$,
donde $x_{ijk}=1$ si el veh\'iculo $k$ viaja directamente del nodo $i$ al nodo
$j$, y $0$ en caso contrario. Aunque la implementaci\'on usa una representaci\'on
evolutiva por permutaciones, el modelo matem\'atico de referencia permite
explicar el objetivo del sistema:

\begin{equation}
\min Z = \sum_{k=1}^{M}\sum_{i=0}^{n}\sum_{j=0}^{n} c_{ij}x_{ijk}
\end{equation}

sujeto a que cada cliente sea visitado una sola vez:

\begin{equation}
\sum_{k=1}^{M}\sum_{i=0}^{n} x_{ijk} = 1, \quad j=1,\ldots,n
\end{equation}

y a que la carga asignada a cada veh\'iculo no exceda su capacidad:

\begin{equation}
\sum_{i=1}^{n} d_i y_{ik} \leq Q, \quad k=1,\ldots,M
\end{equation}

donde $y_{ik}=1$ si el cliente $i$ es atendido por el veh\'iculo $k$. Tambi\'en se
requiere que las rutas salgan del dep\'osito y regresen a \'el. En la aplicaci\'on,
estas restricciones se manejan mediante la separaci\'on de rutas por capacidad y
mediante penalizaciones cuando la soluci\'on requiere m\'as veh\'iculos que los
disponibles.

\subsection{Funci\'on objetivo y penalizaciones}

La implementaci\'on calcula el costo de una soluci\'on con una funci\'on objetivo
penalizada:

\begin{equation}
Z = \sum_{r=1}^{R} Distancia(Ruta_r) + P
\end{equation}

donde $R$ es el n\'umero de rutas generadas y $P$ representa las penalizaciones.
Si una soluci\'on requiere m\'as veh\'iculos que $M$, se agrega una penalizaci\'on
grande por cada veh\'iculo excedente. Si alguna ruta supera la capacidad $Q$,
tambi\'en se agrega una penalizaci\'on proporcional a la violaci\'on. En el c\'odigo
se utiliza la constante:

\begin{equation}
BIG\_PENALTY = 1,000,000
\end{equation}

Este valor no representa una distancia real; su funci\'on es hacer que las
soluciones no factibles sean menos atractivas para el algoritmo. De esta forma,
el algoritmo gen\'etico puede comparar cualquier cromosoma, pero tender\'a a
conservar aquellos que generan rutas factibles y de menor distancia.

\subsection{Restricciones consideradas}

El sistema considera las siguientes restricciones operativas:

\begin{itemize}
    \item Cada cliente debe ser atendido una vez dentro de la permutaci\'on del
    cromosoma.
    \item Cada veh\'iculo inicia su ruta en el dep\'osito y termina en el mismo
    dep\'osito.
    \item La carga total de una ruta no debe superar la capacidad m\'axima $Q$.
    \item El n\'umero de rutas generadas no debe ser mayor que el n\'umero de
    veh\'iculos disponibles $M$.
    \item Un paquete no debe tener exactamente la misma ubicaci\'on que el
    dep\'osito, ya que esto puede causar ambig\"uedad en el c\'alculo y la
    visualizaci\'on.
    \item En modo local, los puntos deben pertenecer a una regi\'on cercana para
    que la red vial descargada por OSMnx sea razonable.
\end{itemize}

Algunas restricciones se validan antes de ejecutar el algoritmo, mientras que
otras se controlan por penalizaci\'on durante la evaluaci\'on de soluciones. Esta
combinaci\'on permite que la aplicaci\'on sea flexible sin dejar de orientar la
b\'usqueda hacia soluciones v\'alidas.

\subsection{Arquitectura de la aplicaci\'on}

La aplicaci\'on sigue una arquitectura modular. Cada archivo principal cumple una
responsabilidad distinta, lo cual facilita el mantenimiento y la explicaci\'on
del sistema. La interfaz no contiene directamente toda la l\'ogica matem\'atica;
en su lugar, llama a funciones especializadas para resolver el VRP y construir
mapas.

\begin{table}
\centering
\caption{M\'odulos principales del sistema.}\label{tab:modulos}
\begin{tabular}{p{3.2cm}p{8.2cm}}
\toprule
M\'odulo & Responsabilidad \\
\midrule
\texttt{desktop\_app.py} & Construye la interfaz gr\'afica, administra eventos,
lee par\'ametros, valida datos, ejecuta la optimizaci\'on y presenta resultados. \\
\texttt{ga\_vrp.py} & Implementa el algoritmo gen\'etico, la funci\'on de
aptitud, la divisi\'on de rutas, la cruza OX y la mutaci\'on. \\
\texttt{map\_utils.py} & Calcula distancias, descarga redes viales, obtiene
rutas en calles reales y calcula distancias aproximadas con Haversine. \\
\texttt{data/clientes.csv} & Contiene el conjunto de datos local usado como
ejemplo de entregas regionales. \\
\texttt{data/clientes\_nacional.csv} & Contiene el conjunto de datos nacional
usado para probar el modo aproximado. \\
\bottomrule
\end{tabular}
\end{table}

\subsection{Descripci\'on de los m\'odulos}

El archivo \texttt{desktop\_app.py} define la clase principal de la aplicaci\'on
de escritorio. Esta clase hereda de \texttt{tk.Tk} y construye los elementos
visuales: panel de configuraci\'on, tabla de paquetes, formulario de captura,
botones de acci\'on, tabla de resumen y ventanas auxiliares. Tambi\'en contiene la
funci\'on \texttt{solve\_vrp}, que conecta la interfaz con el c\'alculo de
distancias, el algoritmo gen\'etico y la generaci\'on del mapa.

El archivo \texttt{ga\_vrp.py} concentra el componente de optimizaci\'on. Incluye
funciones para calcular la carga de una ruta, calcular la distancia recorrida,
dividir una permutaci\'on en rutas seg\'un la capacidad del veh\'iculo, evaluar la
aptitud de una soluci\'on y ejecutar el ciclo evolutivo. Esta separaci\'on es
importante porque permite modificar el algoritmo sin redise\~nar la interfaz.

El archivo \texttt{map\_utils.py} resuelve el componente geogr\'afico. Para rutas
locales, obtiene una red vial con OSMnx y calcula caminos m\'as cortos con
NetworkX. Para rutas nacionales, calcula distancias aproximadas con la f\'ormula
Haversine. Tambi\'en convierte rutas del algoritmo en listas de coordenadas que
pueden dibujarse en Folium.

\subsection{Dise\~no de la interfaz gr\'afica}

La interfaz fue dise\~nada para que el usuario pueda trabajar sin modificar el
c\'odigo. En el panel lateral se configuran los datos globales: cantidad de
veh\'iculos, capacidad m\'axima, coordenadas del dep\'osito, modo de distancia,
tama\~no de poblaci\'on, n\'umero de generaciones y tasa de mutaci\'on. Esta
distribuci\'on concentra los par\'ametros antes de ejecutar la optimizaci\'on.

En el panel principal se encuentra la gesti\'on de env\'ios. El usuario puede
cargar un CSV, guardar los datos actuales, abrir un ejemplo nacional, agregar
filas manualmente, actualizar un destino seleccionado o eliminarlo. Tambi\'en
puede abrir un mapa selector y agregar destinos mediante clics. Esta funci\'on
facilita la captura de coordenadas cuando no se conoce exactamente la latitud y
longitud de un punto.

La parte inferior de la ventana muestra los resultados. La tabla de resumen
indica el veh\'iculo asignado, el color de la ruta, los paquetes transportados,
la carga total, la distancia, el origen, el destino y la secuencia completa de
entrega. Adem\'as, la aplicaci\'on abre el mapa generado en el navegador y permite
abrir una gr\'afica detallada de convergencia.

\begin{figure}
\centering
\fbox{\begin{minipage}{0.90\textwidth}
\centering
Espacio para captura de pantalla de la ventana principal.\\
Debe mostrar el panel de configuraci\'on, la tabla de paquetes y los botones de
ejecuci\'on.
\end{minipage}}
\caption{Ventana principal de la aplicaci\'on de escritorio.}
\label{fig:ventana-principal}
\end{figure}

\begin{figure}
\centering
\fbox{\begin{minipage}{0.90\textwidth}
\centering
Espacio para captura de pantalla del mapa de rutas.\\
Debe mostrar el dep\'osito, los clientes y las rutas coloreadas por veh\'iculo.
\end{minipage}}
\caption{Mapa interactivo generado despu\'es de ejecutar la optimizaci\'on.}
\label{fig:mapa-rutas}
\end{figure}

\begin{figure}
\centering
\fbox{\begin{minipage}{0.90\textwidth}
\centering
Espacio para captura de pantalla de la gr\'afica de convergencia.\\
Debe mostrar la evoluci\'on del mejor valor de la funci\'on objetivo por
generaci\'on.
\end{minipage}}
\caption{Gr\'afica de convergencia del algoritmo gen\'etico.}
\label{fig:convergencia}
\end{figure}

\subsection{Dise\~no del algoritmo gen\'etico}

El algoritmo gen\'etico es el componente central del sistema. Su objetivo es
explorar diferentes \'ordenes de visita de los clientes y conservar aquellas
soluciones que producen menor distancia total penalizada. El dise\~no se basa en
una representaci\'on compacta por permutaciones, operadores adecuados para no
repetir clientes y un mecanismo de elitismo para conservar buenas soluciones.

\subsubsection{Representaci\'on de cromosomas}
Cada cromosoma es una lista ordenada de clientes. El dep\'osito no se incluye en
el cromosoma porque siempre aparece al inicio y al final de cada ruta. Por
ejemplo, si existen cinco clientes, un individuo puede ser:

\begin{equation}
[4, 1, 3, 5, 2]
\end{equation}

Esta lista indica el orden global en que se intentar\'an atender los clientes.
Despu\'es, la funci\'on de divisi\'on de rutas recorre la permutaci\'on y agrupa
clientes hasta que agregar el siguiente exceder\'ia la capacidad $Q$. En ese
momento se cierra la ruta, se regresa al dep\'osito y se abre una nueva ruta.

\subsubsection{Poblaci\'on inicial}
La poblaci\'on inicial se genera aleatoriamente. Cada individuo es una permutaci\'on
distinta de los clientes. Si el tama\~no de poblaci\'on es 50, el sistema crea 50
posibles \'ordenes de visita. Esta diversidad inicial es importante porque da al
algoritmo diferentes zonas de b\'usqueda desde el inicio.

\subsubsection{Funci\'on de aptitud}
La funci\'on de aptitud eval\'ua la calidad de un cromosoma. Primero convierte la
permutaci\'on en rutas, luego calcula la distancia de cada ruta y finalmente suma
las penalizaciones. Como el objetivo es minimizar, una soluci\'on con menor costo
penalizado es considerada mejor.

\subsubsection{Selecci\'on}
La selecci\'on se realiza tomando padres desde un grupo de los mejores individuos
de la poblaci\'on ordenada por aptitud. Esta estrategia favorece soluciones de
buena calidad, pero mantiene cierta variabilidad porque los padres se eligen al
azar dentro de ese grupo superior. En la implementaci\'on se usa un tama\~no de
grupo controlado para evitar que la b\'usqueda dependa de un solo individuo.

\subsubsection{Cruza OX}
La cruza por orden u Order Crossover (OX) es adecuada para cromosomas que son
permutaciones. El operador selecciona un segmento del primer padre y lo copia al
hijo. Luego recorre el segundo padre en orden y coloca los clientes faltantes en
las posiciones vac\'ias. As\'i se preserva una parte del primer padre y parte del
orden relativo del segundo, sin repetir clientes ni perder destinos.

\subsubsection{Mutaci\'on}
La mutaci\'on usada es por intercambio. Con una probabilidad definida por el
usuario, el algoritmo selecciona dos posiciones del cromosoma e intercambia sus
clientes. Este operador introduce variaci\'on y ayuda a escapar de soluciones
estancadas. Una tasa de mutaci\'on demasiado baja puede reducir la exploraci\'on,
mientras que una tasa demasiado alta puede volver inestable la evoluci\'on.

\subsubsection{Elitismo}
El elitismo conserva una porci\'on de los mejores individuos de una generaci\'on y
los copia directamente a la siguiente. Esto evita perder la mejor soluci\'on
encontrada por efectos de cruza o mutaci\'on. En la aplicaci\'on, el tama\~no de
\'elite se calcula con base en el tama\~no de poblaci\'on, manteniendo un m\'inimo y
un m\'aximo razonables.

\subsubsection{Criterio de paro}
El criterio de paro principal es el n\'umero de generaciones configurado por el
usuario. En cada generaci\'on se guarda el mejor costo encontrado, lo que permite
construir la gr\'afica de convergencia. Cuando se alcanza el n\'umero definido de
generaciones, el algoritmo devuelve la mejor soluci\'on, su costo y el historial.

\subsection{Flujo completo de ejecuci\'on}

El flujo de ejecuci\'on inicia cuando el usuario presiona el bot\'on de optimizar.
Primero, la aplicaci\'on valida que existan paquetes y que los par\'ametros sean
num\'ericos. Despu\'es verifica que ning\'un paquete tenga exactamente las mismas
coordenadas que el dep\'osito. A continuaci\'on, construye la lista de coordenadas
y demandas, colocando el dep\'osito como nodo cero.

Si el modo elegido es local, se calcula la separaci\'on m\'axima entre puntos. Si
la distancia es demasiado grande, la aplicaci\'on solicita usar el modo nacional
aproximado. Si la distancia es aceptable, se descarga o genera una red vial
alrededor de las coordenadas. Luego se obtiene la matriz de distancias. Si el
modo elegido es nacional, la matriz se calcula directamente con Haversine.

Con la matriz de distancias lista, se ejecuta el algoritmo gen\'etico. Al final,
la mejor permutaci\'on se divide en rutas, se calcula el resumen por veh\'iculo, se
construye el mapa con Folium y se actualiza la interfaz. La Fig.~\ref{fig:flujo}
resume este procedimiento.

\begin{figure}
\centering
\fbox{\begin{minipage}{0.92\textwidth}
\centering
Inicio $\rightarrow$ leer par\'ametros $\rightarrow$ validar datos
$\rightarrow$ calcular matriz de distancias $\rightarrow$ generar poblaci\'on
$\rightarrow$ evaluar aptitud $\rightarrow$ seleccionar padres $\rightarrow$
cruza OX $\rightarrow$ mutaci\'on $\rightarrow$ elitismo $\rightarrow$
repetir generaciones $\rightarrow$ mostrar rutas, mapa y convergencia.
\end{minipage}}
\caption{Flujo general de ejecuci\'on del sistema.}
\label{fig:flujo}
\end{figure}

\subsection{C\'alculo de distancias con OSMnx y Haversine}

En modo local, el sistema usa OSMnx para obtener una red de calles manejables
alrededor de los puntos. Cada coordenada se asocia al nodo m\'as cercano de la
red. Posteriormente, NetworkX calcula las distancias m\'as cortas entre nodos
usando el peso de longitud de las calles. Para mejorar la aproximaci\'on, el
sistema agrega la distancia entre la coordenada exacta y el nodo vial m\'as
cercano.

En modo nacional, el sistema usa la f\'ormula Haversine, que estima la distancia
entre dos puntos sobre la superficie terrestre a partir de latitud y longitud.
La expresi\'on general es:

\begin{equation}
a = \sin^2\left(\frac{\Delta \varphi}{2}\right) +
\cos(\varphi_1)\cos(\varphi_2)\sin^2\left(\frac{\Delta \lambda}{2}\right)
\end{equation}

\begin{equation}
d = 2R\arcsin(\sqrt{a})
\end{equation}

donde $R$ es el radio de la Tierra, $\varphi$ representa latitud y $\lambda$
representa longitud. Debido a que una carretera no suele seguir una l\'inea
recta, la distancia se multiplica por un factor de carretera. En la aplicaci\'on
se utiliza un factor de 1.25 para obtener una estimaci\'on m\'as realista que la
distancia geod\'esica pura.

\subsection{Diagramas UML y diagramas del sistema}

El primer diagrama representa los casos de uso principales. El actor central es
el usuario, quien puede configurar veh\'iculos, cargar clientes, editar la tabla,
seleccionar destinos en el mapa, ejecutar la optimizaci\'on y consultar
resultados.

\begin{figure}
\centering
\fbox{\begin{minipage}{0.92\textwidth}
\centering
Usuario: configurar veh\'iculos; cargar CSV; editar paquetes; seleccionar
destinos en mapa; ejecutar optimizaci\'on; consultar tabla de rutas; abrir mapa;
abrir gr\'afica de convergencia.
\end{minipage}}
\caption{Diagrama textual de casos de uso del sistema.}
\label{fig:casos-uso}
\end{figure}

El segundo diagrama resume las clases y componentes l\'ogicos. La clase
\texttt{VRPDesktopApp} coordina la interfaz. El m\'odulo \texttt{ga\_vrp} ofrece
funciones de optimizaci\'on y el m\'odulo \texttt{map\_utils} ofrece funciones de
distancia y mapas. La relaci\'on principal es de uso: la interfaz invoca a los
m\'odulos de c\'alculo cuando el usuario ejecuta la optimizaci\'on.

\begin{figure}
\centering
\fbox{\begin{minipage}{0.92\textwidth}
\centering
\texttt{VRPDesktopApp} $\rightarrow$ \texttt{solve\_vrp}
$\rightarrow$ \texttt{map\_utils} para distancias y rutas.\\
\texttt{solve\_vrp} $\rightarrow$ \texttt{ga\_vrp} para algoritmo gen\'etico.\\
\texttt{solve\_vrp} $\rightarrow$ \texttt{build\_map} para visualizaci\'on.
\end{minipage}}
\caption{Diagrama textual de componentes del sistema.}
\label{fig:componentes}
\end{figure}

\section{Experimentaci\'on}

\subsection{Ambiente de desarrollo}

La aplicaci\'on fue desarrollada en Python sobre un entorno virtual local. El
proyecto se organiz\'o en una carpeta de trabajo con archivos fuente, datos de
prueba, cach\'e de mapas y carpeta de salida para los mapas generados. El
ambiente de desarrollo usado para las pruebas fue una computadora personal con
sistema operativo Windows, ejecutando la aplicaci\'on desde PowerShell mediante
el archivo \texttt{run\_app.bat} o directamente con \texttt{python
desktop\_app.py}.

\begin{table}
\centering
\caption{Ambiente de software utilizado.}\label{tab:software}
\begin{tabular}{p{4cm}p{7.4cm}}
\toprule
Elemento & Descripci\'on \\
\midrule
Lenguaje & Python \\
Interfaz gr\'afica & Tkinter y ttk \\
Mapas interactivos & Folium \\
Red vial & OSMnx y OpenStreetMap \\
Grafos y caminos m\'as cortos & NetworkX \\
Datos tabulares & CSV y Pandas en la versi\'on Streamlit \\
Visualizaci\'on auxiliar & Navegador web para archivos HTML generados \\
\bottomrule
\end{tabular}
\end{table}

\subsection{Bibliotecas utilizadas}

Las bibliotecas principales son \texttt{folium}, \texttt{osmnx},
\texttt{networkx}, \texttt{pandas}, \texttt{numpy}, \texttt{streamlit} y
\texttt{streamlit-folium}. Aunque la aplicaci\'on principal es de escritorio, se
conserva una versi\'on Streamlit anterior. Folium se usa para producir mapas en
HTML; OSMnx descarga o construye redes viales; NetworkX calcula caminos m\'as
cortos; y Tkinter permite construir la interfaz local.

\subsection{Configuraci\'on de par\'ametros}

Los par\'ametros del algoritmo gen\'etico se definieron desde la interfaz. Para
las pruebas base se consider\'o una poblaci\'on de 50 individuos, 100 generaciones
y una tasa de mutaci\'on de 10\%. Estos valores permiten observar la convergencia
sin hacer demasiado lenta la ejecuci\'on. La interfaz permite aumentar la
poblaci\'on hasta 300, las generaciones hasta 700 y ajustar la mutaci\'on entre
0\% y 100\%.

\begin{table}
\centering
\caption{Par\'ametros base del algoritmo gen\'etico.}\label{tab:parametros}
\begin{tabular}{ll}
\toprule
Par\'ametro & Valor base \\
\midrule
Tama\~no de poblaci\'on & 50 individuos \\
N\'umero de generaciones & 100 generaciones \\
Tasa de mutaci\'on & 0.10 \\
Operador de cruza & OX \\
Operador de mutaci\'on & Intercambio de dos clientes \\
Criterio de paro & N\'umero fijo de generaciones \\
\bottomrule
\end{tabular}
\end{table}

\subsection{Conjuntos de datos empleados}

Para probar la aplicaci\'on se prepararon dos conjuntos de datos. El primero
corresponde a entregas locales en el Istmo de Tehuantepec, Oaxaca, con ocho
destinos y demandas entre 3 y 6 unidades. El segundo conjunto incluye destinos
nacionales como Guadalajara, Tijuana, Canc\'un, M\'erida, Ciudad de M\'exico,
Monterrey, Puebla y Villahermosa. El conjunto local se utiliza con el modo de
calles reales, mientras que el nacional se utiliza con el modo aproximado.

\begin{table}
\centering
\caption{Ejemplo de datos locales utilizados en la experimentaci\'on.}
\begin{tabular}{llll}
\toprule
Destino & Latitud & Longitud & Demanda \\
\midrule
El Espinal & 16.485600 & -95.040300 & 5 \\
Asunci\'on Ixtaltepec & 16.503900 & -95.061900 & 4 \\
Ciudad Ixtepec & 16.562000 & -95.103000 & 6 \\
La Ventosa & 16.551900 & -94.947600 & 5 \\
Uni\'on Hidalgo & 16.472800 & -94.829900 & 4 \\
Santa Mar\'ia Xadani & 16.361600 & -95.018900 & 3 \\
Santo Domingo Ingenio & 16.586900 & -94.761900 & 6 \\
San Blas Atempa & 16.326900 & -95.225600 & 4 \\
\bottomrule
\end{tabular}
\end{table}

\subsection{Procedimiento experimental}

La experimentaci\'on sigui\'o un procedimiento repetible. Primero se carg\'o el
conjunto de datos local y se revis\'o que las demandas fueran positivas. Despu\'es
se definieron el dep\'osito, el n\'umero de veh\'iculos y la capacidad. Enseguida
se ejecut\'o el algoritmo con los par\'ametros base y se registraron las rutas
obtenidas. Finalmente, se analiz\'o el mapa y la gr\'afica de convergencia. El
mismo proceso se repiti\'o con el conjunto nacional, cambiando el modo de
distancia.

\begin{enumerate}
    \item Cargar el archivo CSV de clientes locales.
    \item Definir el dep\'osito inicial mediante latitud y longitud.
    \item Configurar la cantidad de veh\'iculos disponibles y la capacidad
    m\'axima por veh\'iculo.
    \item Seleccionar el modo de distancia local con calles OSMnx para el caso
    regional.
    \item Ejecutar el algoritmo gen\'etico con una poblaci\'on inicial, un n\'umero
    de generaciones y una tasa de mutaci\'on configurables.
    \item Revisar la tabla de rutas, el mapa generado y la gr\'afica de
    convergencia.
    \item Repetir la prueba con el conjunto nacional usando el modo aproximado.
\end{enumerate}

\subsection{Variables de entrada y salida}

Las variables de entrada principales son el n\'umero de veh\'iculos $M$, la
capacidad $Q$, la lista de clientes, las demandas, el dep\'osito y los par\'ametros
del algoritmo gen\'etico. Las salidas son la mejor permutaci\'on encontrada, las
rutas separadas por veh\'iculo, la distancia total penalizada, la factibilidad de
la soluci\'on, el mapa interactivo y el historial de convergencia.

\begin{table}
\centering
\caption{Variables principales de entrada y salida.}\label{tab:variables}
\begin{tabular}{p{3.5cm}p{8cm}}
\toprule
Tipo & Variables \\
\midrule
Entrada & Veh\'iculos disponibles, capacidad, dep\'osito, destinos, demandas,
poblaci\'on, generaciones, tasa de mutaci\'on y modo de distancia. \\
Salida & Cromosoma ganador, rutas por veh\'iculo, carga por ruta, distancia,
factibilidad, mapa HTML y gr\'afica de convergencia. \\
\bottomrule
\end{tabular}
\end{table}

\subsection{Resultados obtenidos}

Los resultados se presentan como tablas de resumen por veh\'iculo. Debido a que
el algoritmo gen\'etico incluye aleatoriedad, los valores exactos pueden cambiar
entre ejecuciones. Por ello, la evaluaci\'on se enfoca en verificar que la
soluci\'on generada cumpla la estructura esperada: rutas que salen del dep\'osito,
atienden paquetes, regresan al dep\'osito, respetan capacidad y no exceden el
n\'umero de veh\'iculos cuando existe una soluci\'on factible.

\begin{table}
\centering
\caption{Formato de resultados por veh\'iculo generado por la aplicaci\'on.}
\label{tab:resultados}
\begin{tabular}{llll}
\toprule
Veh\'iculo & Paquetes & Carga total & Distancia \\
\midrule
1 & P1, P2, P6 & Seg\'un demandas & Calculada por matriz \\
2 & P3, P4 & Seg\'un demandas & Calculada por matriz \\
3 & P5, P7, P8 & Seg\'un demandas & Calculada por matriz \\
\bottomrule
\end{tabular}
\end{table}

La tabla anterior representa el tipo de salida obtenida. En la aplicaci\'on real,
cada fila tambi\'en muestra el color de la ruta, el origen, el destino y la
secuencia completa. Esto permite verificar r\'apidamente si un veh\'iculo tiene
demasiados paquetes o si la ruta asignada es coherente con el mapa.

\subsection{Gr\'aficas de convergencia}

La gr\'afica de convergencia muestra el mejor valor de la funci\'on objetivo en
cada generaci\'on. En una ejecuci\'on t\'ipica, el costo disminuye durante las
primeras generaciones porque la poblaci\'on inicial contiene soluciones
aleatorias y el algoritmo encuentra combinaciones mejores. Despu\'es, la curva
puede estabilizarse cuando las mejoras se vuelven menos frecuentes.

Esta gr\'afica permite tomar decisiones sobre los par\'ametros. Si la curva sigue
mejorando al final de las generaciones, conviene aumentar el n\'umero de
generaciones. Si se estanca muy pronto, se puede aumentar la mutaci\'on o la
poblaci\'on para explorar m\'as soluciones. Si la curva es muy irregular, una tasa
de mutaci\'on demasiado alta podr\'ia estar afectando la conservaci\'on de buenas
soluciones.

\subsection{Mapas generados}

El mapa generado por la aplicaci\'on representa visualmente la soluci\'on. El
dep\'osito aparece con un marcador especial y cada veh\'iculo se dibuja con un
color diferente. En modo local, las rutas siguen la red vial calculada con
OSMnx. En modo nacional, las rutas se representan como l\'ineas directas entre
puntos, ya que se trata de una aproximaci\'on y no de una red carretera completa.

El mapa cumple dos funciones. Primero, permite validar visualmente la soluci\'on:
si una ruta parece incoherente, el usuario puede revisar los datos o los
par\'ametros. Segundo, comunica el resultado de forma m\'as clara que una tabla,
especialmente para usuarios que desean identificar zonas geogr\'aficas de reparto.

\subsection{Discusi\'on de resultados}

Los resultados muestran que el algoritmo gen\'etico es adecuado para construir
soluciones aproximadas del VRP en el contexto del proyecto. Su principal ventaja
es la flexibilidad: puede trabajar con diferentes cantidades de clientes,
veh\'iculos y capacidades sin cambiar el modelo completo. Adem\'as, la gr\'afica de
convergencia permite observar el comportamiento del proceso de b\'usqueda.

La integraci\'on con mapas tambi\'en es una ventaja importante. En modo local, el
uso de calles reales aporta una interpretaci\'on m\'as cercana a la operaci\'on
log\'istica. En modo nacional, la aproximaci\'on por Haversine permite hacer
pruebas con puntos lejanos sin descargar redes demasiado grandes.

Entre las limitaciones observadas se encuentra la dependencia de internet para
descargar datos de OpenStreetMap en modo local. Tambi\'en se observ\'o que, si los
puntos est\'an muy separados, la descarga de una red vial local deja de ser
pr\'actica. Otra limitaci\'on es que el algoritmo gen\'etico no garantiza el \'optimo
global, por lo que diferentes ejecuciones pueden producir rutas ligeramente
distintas.

\section{Conclusiones}

El proyecto cumpli\'o el objetivo general de desarrollar una aplicaci\'on capaz de
optimizar rutas de reparto mediante un algoritmo gen\'etico. La soluci\'on permite
configurar una flota, cargar clientes, definir demandas, calcular distancias,
ejecutar la optimizaci\'on y visualizar los resultados en tablas, mapas y
gr\'aficas. Esto demuestra que el sistema no se limita al c\'alculo interno, sino
que ofrece una experiencia completa de uso para analizar escenarios de reparto.

Los objetivos espec\'ificos tambi\'en se cumplieron. El problema fue modelado como
un VRP con dep\'osito, clientes, demandas y capacidad. Se implement\'o un algoritmo
gen\'etico con cromosomas por permutaci\'on, poblaci\'on inicial aleatoria, funci\'on
de aptitud penalizada, cruza OX, mutaci\'on por intercambio y elitismo. Adem\'as,
se dise\~n\'o una interfaz gr\'afica que permite administrar los datos y se
incluyeron mapas interactivos para interpretar las rutas resultantes.

El algoritmo gen\'etico aport\'o una estrategia de b\'usqueda flexible para un
problema combinatorio complejo. Su uso permiti\'o generar soluciones aproximadas
sin revisar todas las combinaciones posibles. Tambi\'en permiti\'o visualizar la
evoluci\'on de la soluci\'on mediante la gr\'afica de convergencia. Este aspecto es
relevante porque ayuda a comprender que la optimizaci\'on no ocurre en un solo
paso, sino mediante una mejora progresiva de la poblaci\'on.

La aplicaci\'on desarrollada presenta varias ventajas. La primera es que integra
en una sola herramienta la captura de datos, la optimizaci\'on y la visualizaci\'on.
La segunda es que ofrece dos modos de distancia: uno local con calles reales y
otro nacional aproximado. La tercera es que permite modificar par\'ametros del
algoritmo sin cambiar el c\'odigo, lo cual facilita la experimentaci\'on. La cuarta
es que muestra resultados comprensibles para el usuario, no solo valores
num\'ericos.

Tambi\'en se identificaron limitaciones. El modo local depende de datos de
OpenStreetMap y puede requerir conexi\'on a internet. El modo nacional no calcula
rutas carreteras exactas, sino distancias aproximadas. El algoritmo gen\'etico,
por su naturaleza, no garantiza encontrar el \'optimo global. Adem\'as, el modelo
actual considera veh\'iculos homog\'eneos, un solo dep\'osito y no incluye ventanas
de tiempo, horarios de servicio, tr\'afico, prioridad de clientes ni costos
variables.

Como trabajo futuro se propone ampliar el modelo con ventanas de tiempo, m\'ultiples
dep\'ositos, veh\'iculos con diferentes capacidades, tiempos de descarga, costos de
combustible y restricciones de horario. Tambi\'en podr\'ia integrarse una API de
mapas para obtener rutas y tiempos de viaje m\'as precisos. Otra mejora relevante
ser\'ia comparar el algoritmo gen\'etico con otras t\'ecnicas de optimizaci\'on,
como b\'usqueda tab\'u, recocido simulado, colonias de hormigas o algoritmos
exactos para instancias peque\~nas. Finalmente, se podr\'ia automatizar la
exportaci\'on de reportes en PDF para documentar cada ejecuci\'on.

\section*{Evaluaci\'on Ordinaria}

El reporte cumple con los apartados solicitados para la Evaluaci\'on Ordinaria:
Introducci\'on, Objetivos, Dise\~no, Experimentaci\'on y Conclusiones. Adem\'as,
presenta la aplicaci\'on bajo una estructura compatible con la plantilla LNCS de
Springer, usando resumen, palabras clave, secciones formales, ecuaci\'on de
funci\'on objetivo y tabla de datos experimentales.

\begin{credits}
\subsubsection{\ackname} Se agradece el uso de herramientas de c\'odigo abierto
como Python, Tkinter, Folium, OSMnx y NetworkX, que permitieron construir y
probar la aplicaci\'on descrita en este reporte.

\subsubsection{\discintname}
El autor declara que no existen intereses en competencia relacionados con el
contenido de este reporte.
\end{credits}

% ---- Bibliography ----
%
% BibTeX users should specify bibliography style 'splncs04'.
% References will then be sorted and formatted in the correct style.
%
% \bibliographystyle{splncs04}
% \bibliography{mybibliography}
%
\begin{thebibliography}{8}

\bibitem{toth2014}
Toth, P., Vigo, D.: Vehicle Routing: Problems, Methods, and Applications.
Society for Industrial and Applied Mathematics, Philadelphia (2014)

\bibitem{goldberg1989}
Goldberg, D.E.: Genetic Algorithms in Search, Optimization, and Machine
Learning. Addison-Wesley, Boston (1989)

\bibitem{osmnx}
Boeing, G.: OSMnx: New methods for acquiring, constructing, analyzing, and
visualizing complex street networks. Computers, Environment and Urban Systems
65, 126--139 (2017)

\bibitem{folium}
Folium contributors: Folium documentation. \url{https://python-visualization.github.io/folium/}

\bibitem{networkx}
Hagberg, A.A., Schult, D.A., Swart, P.J.: Exploring network structure,
dynamics, and function using NetworkX. In: Proceedings of the 7th Python in
Science Conference, pp. 11--15 (2008)

\end{thebibliography}

\end{document}
